% !TEX root = ../main.tex
\PlanPage{协作维护、版本治理与研究归档}{吕洪宇}{41}
\noindent\textbf{项目协作的目标不是简单记录谁完成了哪项任务，而是保证每一次更新都能回答“改动了什么、依据是什么、影响了哪份研究报告”。}团队以Issue拆分可验收任务，以分支和拉取请求承载代码、数据与文档变更，以统一JSON契约连接模块，以运行记录和证据账本追溯结果。计划书由吕洪宇与武桐宇负责撰写，研发模块的贡献则按团队文档中的实际分工保留；两类责任在材料中分别说明\cite{lhy-team,lhy-data-standard}。

\begin{center}\small 表41-1\quad 代码、数据与计划书的共同归档规则\par\vspace{1mm}
\begin{tabularx}{\linewidth}{p{0.18\linewidth}p{0.34\linewidth}X}
\toprule
归档对象 & 每次变更应记录的内容 & 审阅与定位入口\\
\midrule
代码与接口 & 修改理由、字段兼容性、正常/异常测试、影响模块与提交号 & Issue、PR、提交记录与契约评审；跨模块字段变化需关联调用方。\\
图谱、证据和评分数据 & 来源、日期、口径、数据版本、核验状态、责任人与适用范围 & 数据文件、来源复核文档和证据账本；敏感凭据与受限原文不进入仓库。\\
运行与演示结果 & 事件ID、数据版本、\texttt{run\_id}、报告ID、生成时间、异常与截图模式 & 本地运行记录、演示脚本和截图清单；截图可回溯到实现提交。\\
比赛计划书 & 页位、图表来源、引用键、编译产物、逐页检查与交叉审阅意见 & 写作分支、PDF预览、证据账本、自检表及合并记录。\\
\bottomrule
\end{tabularx}\end{center}

在研究结果层面，最小归档单元是一份能够重演的事件报告：它至少包含输入事件、数据版本、代码提交、报告输出、路径和公司记录、证据来源及核验状态。若研究结论在后续版本中变化，团队可比较两个报告的事件、图谱、参数、证据和代码差异，进而定位变化的直接原因。这个方法避免把系统迭代笼统归因为“模型变化”，使每次修订都落在可检查的材料上。

版本治理同样服务于比赛展示。评审看到的每一张系统截图、每一个表格数值和每一段功能说明，都应当对应固定的数据状态与来源；当项目进入新阶段时，团队以新版本补充证据，而不是回写或混淆旧版本的报告口径。由此，协作过程成为研究可信度的一部分。
